% !PRINTED 474
\weberpage{474}

%<<CONT
\[
\pi_a = \begin{pmatrix} 0, & 1 & \ldots & m-1 \\ a_0, & a_1 & \ldots & a_{m-1} \end{pmatrix}, \quad
\pi_b = \begin{pmatrix} 0, & 1 & \ldots & m-1 \\ b_0, & b_1 & \ldots & b_{m-1} \end{pmatrix}
\]
are any two permutations, then by (6) we can also write the second
thus:
\[
\pi_b = \begin{pmatrix} a_0, & a_1 & \ldots & a_{m-1} \\ b_{a_0}, & b_{a_1} & \ldots & b_{a_{m-1}} \end{pmatrix},
\]
and if we therefore carry out both one after the other, first $\pi_a$,
then $\pi_b$, there arises a third permutation, which we call composed
(componirt) from $\pi_a$ and $\pi_b$, and which we represent as the
symbolic product of $\pi_a$ and $\pi_b$:

\weq{7}{\pi_a \pi_b = \begin{pmatrix} 0, & 1 & \ldots & m-1 \\ b_{a_0}, & b_{a_1} & \ldots & b_{a_{m-1}} \end{pmatrix}.}

In this composition the factors usually cannot be interchanged,
since $\pi_a \pi_b$ does not in general agree with
$\pi_b \pi_a$. The following symbolic equation expresses the associative
principle.

\weq{8}{(\pi_a \pi_b)\pi_c = \pi_a(\pi_b \pi_c) = \pi_a \pi_b \pi_c}

One convinces oneself of this by the following simple
consideration.

If an index $p$ is changed by $\pi_a$ into $q$, the resulting index is changed
by $\pi_b$ into $r$, and $r$ by $\pi_c$ into $s$, then $p$ passes by
$\pi_a \pi_b$ into $r$, and $q$ by $\pi_b \pi_c$ into $s$; hence
$p$ is changed into $s$ both by $(\pi_a \pi_b)\pi_c$ and by
$\pi_a(\pi_b \pi_c)$. But formula (8) says nothing else.

To every permutation $\pi_a$ there is a definite
\gesperrt{opposite} (inverse) permutation, which consists in this,
that the change produced by $\pi_a$ is undone again, and which
therefore, composed with $\pi_a$, produces the identical
permutation. We denote this permutation by
$\pi_a^{-1}$, and we have, if $\pi_a$ has the expression (5),

\weq{9}{\pi_a^{-1} = \begin{pmatrix} a_0, & a_1 & \ldots & a_{m-1} \\ 0, & 1 & \ldots & m-1 \end{pmatrix}.}

For by the rule of composition in (7) we have
\[
\pi_a \pi_a^{-1} = \begin{pmatrix} 0, & 1 & \ldots & m-1 \\ 0, & 1 & \ldots & m-1 \end{pmatrix}.
\]
But in the same way we also have
\[
\pi_a^{-1}\pi_a = \begin{pmatrix} a_0, & a_1 & \ldots & a_{m-1} \\ a_0, & a_1 & \ldots & a_{m-1} \end{pmatrix},
\]
%CONT>>
